\section{Introduction}\label{sec:intro}

Data agents convert natural-language questions into executable
data pipelines. They discover sources, inspect schemas, write SQL
or Python programs, repair execution errors, and report computed
answers with supporting evidence
\citep{kramabench2025,dabstep2025,lakeqa2026}.
This workflow changes the object of verification. A reported value
may occur in no source file because it is produced by filtering,
joining, and aggregation. Auditing a computed answer therefore
requires evidence about the execution that produced it, rather than
only evidence that contains related text or values.

A KramaBench episode illustrates the first failure. The question
requests the average number of identity-theft reports per 100,000
population across metropolitan areas. The agent cites 17 files and
states 10,724.14. Its final program executes, yet replay under the
cited-file view produces no value comparable to the stated answer.
The citations are relevant, but the recorded computation does not
reproduce the claim. A second failure arises even when replay
succeeds. A program can average the wrong years, select the wrong
subset, or compute a sum when the question requests a mean. The
reported value is then reproducible, while the implemented operation
is semantically incorrect.

These failures expose three relations that a post-execution audit
must establish. First, \emph{evidence authorization} relates the
sources declared by the agent to the sources actually read by the
host program. Second, \emph{execution reproducibility} relates the
final program and its authorized evidence view to the stated answer.
Third, \emph{semantic correctness} relates the final program to the
task specification. The first two relations admit deterministic
trace and execution witnesses. The third requires a task-specific
specification, such as reference SQL, executable constraints,
database variants, metamorphic relations, or human judgement.
A complete audit record must preserve these relations separately:
a reproducible answer may still implement the wrong task, while a
semantically appropriate program may lack sufficient execution
evidence for automatic delivery.

Existing methods address adjacent parts of this chain. Citation
metrics evaluate whether generated text is supported by retrieved
passages \citep{gao2023alce,min2023factscore}. Database provenance
represents how a known query output depends on input tuples
\citep{green2007provenance,cheney2009provenance,
amsterdamer2011aggregate}, while script-provenance systems capture
runtime dependencies in controlled environments
\citep{murta2014noworkflow,namaki2020vamsa}. Agent benchmarks and
model-based judges evaluate final utility or interaction risk
\citep{ruan2024toolemu,debenedetti2024agentdojo,kramabench2025,
sana2026}. These mechanisms do not provide a deterministic
per-answer witness that jointly connects the final program, stated
answer, declared evidence, and observed data reads of a completed
data-agent episode.

We introduce \ledger, a deterministic post-execution audit layer
for computed data-agent answers. \ledger first records an explicit
trace contract containing the final program, stated answer, declared
sources, observed host reads, execution outcome, and reached
operands. A canonical resolver and read-reconciliation step check
evidence authorization. Citation-restricted replay then executes the
final program inside the resolved evidence view and compares its
output with the stated answer, producing a replay certificate or a
localized failure. When the original runtime is unavailable,
operand-constrained closure provides a bounded derivation witness or
returns unknown. A separate task-specific semantic interface checks
the question--program relation when an executable specification is
available. The provenance and semantic statuses remain independent
and are jointly consumed by a governance policy for delivery,
abstention, or escalation.

We make four contributions.
\begin{enumerate}
    \setlength{\itemsep}{0pt}
    \setlength{\parskip}{0pt}

    \item \textbf{Problem formulation.}
    We formulate post-execution auditing for computed data-agent
    answers as three relations: evidence authorization, execution
    reproducibility, and semantic correctness. The formulation
    includes an explicit trace contract, trusted computing base,
    failure semantics, and an unknown outcome for insufficient
    evidence.

    \item \textbf{Theory.}
    We establish the relative soundness of citation-restricted
    replay under the stated trust assumptions, prove completeness
    of the bounded operand fallback within its operation and depth
    limits, and formalize why execution-provenance consistency does
    not imply semantic correctness.

    \item \textbf{Method.}
    We develop \ledger, which combines canonical evidence
    resolution, declared-versus-observed read reconciliation,
    citation-restricted replay, type-aware answer comparison,
    operand-constrained closure, and an independent task-specific
    semantic-validation interface. Every audit outcome carries a
    machine-readable witness or an explicit unknown status.

    \item \textbf{Evaluation.}
    We evaluate the three relations across increasingly realistic
    evidence regimes. On WikiSQL, the semantic interface detects
    controlled program changes that preserve the original execution
    result. On ConvFinQA and MultiHiertt, the provenance audit
    recovers 89.6--93.3\% of reference derivations while rejecting
    96.1--97.2\% of cross-example unsupported answers. On
    human-labelled and scored live data-agent episodes, replay-based
    selection raises the proportion of delivered answers judged
    supported and improves mean delivered score at partial coverage.
\end{enumerate}
